\begin{table}[H]
\centering
\caption{Standardized Effect Sizes for Main Outcomes}
\label{tab:sde}
\begin{adjustbox}{max width=\textwidth}
\begin{tabular}{llccccc}
\toprule
Outcome & Specification & $\hat{\beta}$ & SD($Y$) & SDE & SE(SDE) & Classification \\
\midrule
\multicolumn{7}{l}{\textit{Panel A: Pooled}} \\
Export (MW) & 2021 reform (6h$\to$4h) & 55.7 & 1446.7 & 0.0385 & 0.1014 & Small positive \\
Export (MW) & 2024 reform (4h$\to$3h) & -185.8 & 1161.9 & -0.1599 & 0.1360 & Large negative \\
Export (MW) & Pooled (both reforms) & -74.2 & 1385.2 & -0.0536 & 0.0632 & Moderate negative \\
\midrule
\multicolumn{7}{l}{\textit{Panel B: Heterogeneous (by interconnector capacity)}} \\
Export (MW) & High-capacity neighbors & 162.4 & 1636.6 & 0.0992 & 0.1687 & Moderate positive \\
Export (MW) & Low-capacity neighbors & -9.5 & 554.3 & -0.0172 & 0.1376 & Small negative \\
\bottomrule
\end{tabular}
\end{adjustbox}
\par\vspace{0.5em}
\begin{minipage}{\textwidth}\sloppy
\footnotesize\textit{Notes:} \textbf{Country:} Germany and 11 European neighbors (Austria, Belgium, Czechia, Denmark, France, Luxembourg, Netherlands, Norway, Poland, Sweden, Switzerland). \textbf{Research question:} Whether tightening Germany's EEG \S 51 subsidy clawback threshold for renewable generators reduces bilateral electricity exports to neighboring countries during negative-price episodes. \textbf{Policy mechanism:} The EEG clawback suspends market premium payments to renewable generators above 400 kW when the day-ahead electricity price is negative for N or more consecutive hours; threshold tightened from 6h to 4h (January 2021) and from 4h to 3h (January 2024), creating financial incentives for generators to curtail output during extended negative-price periods. \textbf{Outcome definition:} Mean hourly bilateral electricity export from Germany to each neighbor (MW) during a negative-price episode, where positive values indicate power flowing from Germany to the neighbor. \textbf{Treatment:} Binary; equals one for episodes whose duration falls in the newly clawback-eligible window created by each threshold tightening. \textbf{Data:} Fraunhofer ISE Energy-Charts API, hourly bilateral cross-border flows and DE-LU day-ahead prices, 2019--2025, episode-neighbor level, 1,334--1,763 observations. \textbf{Method:} Difference-in-differences comparing newly clawback-eligible episodes (treated) to shorter episodes (control) before and after each threshold change, with neighbor and year-month fixed effects, standard errors clustered at the year-month level. \textbf{Sample:} All consecutive negative-price episodes of 1--5 hours in the DE-LU day-ahead market (excluding episodes $\geq$6 hours that were always clawback-eligible), matched to bilateral flows for 11 neighboring countries. SDE $= \hat{\beta} / \text{SD}(Y)$ where SD($Y$) is the pre-treatment standard deviation. Classification refers to magnitude, not statistical significance: Large ($|$SDE$| > 0.15$), Moderate ($0.05$--$0.15$), Small ($0.005$--$0.05$), Null ($< 0.005$). Classification labels refer to the magnitude of the standardized point estimate, not to statistical significance. ``Null'' denotes a near-zero effect size ($|$SDE$| < 0.005$), not a failure to reject a null hypothesis.
\end{minipage}
\end{table}
